\section{Literature review}\label{sec:literature}

\subsection{LCZs as climate-form vocabulary and weak-label infrastructure}

LCZs were designed as physically interpretable classes for near-surface urban
climate comparison \citep{stewart2012LocalClimate}. Early remote-sensing
classification, monitoring-network design, and worldwide mapping established
the scheme as both a measurement vocabulary and a scalable mapping task
\citep{bechtel2012ClassifiLocal,lelovics2014DesignUrban,
bechtel2015MappingLocal}. WUDAPT extended that logic into a shared urban
modelling infrastructure, but comparisons with detailed GIS canopy parameters
also showed that class-based descriptions cannot reproduce every local
morphological quantity \citep{ching2018WudaptUrban,
hammerberg2018ImplicatEmployin,zonato2020EvaluatiPerforma}.

Classification performance depends on sensor, city, training data, object or
pixel support, and feature fusion. Random forests, convolutional networks,
residual networks, multilevel fusion, and multisource seasonal imagery have all
improved LCZ mapping \citep{qiu2018FeatureImportan,yoo2019ComparisBetween,
qiu2020MultilevFeature,yoo2020ImprovinLocal,wang2023LocalClimate}. Yet high
map accuracy does not remove uncertainty caused by incomplete buildings,
different urban forms, or transfer between climatic contexts
\citep{wang2018AssessinLocal,yoo2020ImprovinLocal,
yan2022ComparinObject}. Global mapping is therefore best treated as a
consistent prior whose local meaning must be checked, rather than as a complete
substitute for local labels \citep{demuzere2022GlobalLocal}.

Thermal separability is similarly conditional. Numerical, satellite, and
station evidence shows useful LCZ contrasts, but also within-class variation
and context sensitivity \citep{geletic2016LandSurface,skarbit2017EmployinUrban,
kwok2019WellDoes,yang2020InvestigDiversit}. Frontal area, vegetation, and
scenario definition can change class-level temperature patterns
\citep{li2022VariabilLand,zhou2023IdentifyEffects,
deng2024SensitivLocal}. Coastal-city risk research adds a transferability
warning: an LCZ may organize exposure, but it cannot on its own represent
vulnerability, accessibility, or local design opportunity
\citep{ma2023InvestigUrban,zhang2025ApplicabLocal}. These findings motivate our
weak-label terminology, confidence threshold, mixed-cell flag, and manual-audit
gate.

\subsection{Urban morphometrics, figure-ground form, and fabric typology}

Urban-climate mechanisms predate LCZ classification. Comparative studies linked
urban heat to form and thermal properties, while view-factor research formalized
how surrounding geometry controls radiative exposure \citep{emmanuel2007UrbanHeat,
chen2010ViewFactor}. Sky-view-factor footprints demonstrate that even a
seemingly local 3D metric has a spatial support that must be specified
\citep{middel2018ViewFactor}. Building density, compactness, height,
fragmentation, imperviousness, and vegetation consequently need to be measured
as a system rather than as interchangeable proxies.

Empirical studies have connected LST to building form, landscape configuration,
and multisource urban structure using regression and machine learning
\citep{ferreira2019ExplorinRelation,huang2019InvestigEffects,
sun2019QuantifyEffects,zhang2019EvaluatiEffect}. Later work strengthened the
3D component by evaluating vertical morphology, seasonal disparity, and
scale-dependent landscape patterns \citep{alexander2021InfluencProporti,
yang2021AssessinEffects,chen2022SeasonalDisparat,
han2022UnderstaSeasonal,wu2022QuantifyInfluenc,wu2022ScaleEffect}. The
consistent lesson is not that one feature universally dominates, but that
horizontal density, vertical structure, open space, and vegetation interact
across spatial supports.

Recent studies increasingly compare grid, block, functional-zone, and LCZ
supports. Grid--block differences have been demonstrated directly
\citep{jeon2023ImpactsUrban}; cumulative morphology, satellite/in-situ
comparison, and functional-zone analysis further show that aggregation choices
shape interpretation \citep{park2023QuantifyCumulati,puche2023InsightsInto,
du2024DoesUrban,lin2024DoesUrban}. Studies of stationarity, threshold effects, and urban-rural
or climate-zone differences caution against universal linear prescriptions
\citep{yuan2024EffectsUrban,zhao2024MechanisStationa,
chen2025ImpactsThreshol,debray2025UniversaPatterns}. This literature supports
our decision to report class profiles and transfer diagnostics rather than a
single pooled coefficient.

Three-dimensional green structure is equally important. Canopy amount,
vertical pattern, building--green interaction, and nature-based-solution scale
can modify cooling potential \citep{bai2025ImpactBuilding,
hu2025ComparinGreen,wei2025UrbanCooling}. Surface and canopy temperatures may
also respond differently to 2D and 3D morphology and to analytical scale
\citep{wu2025ContrastUrban,yu2026UrbanMorpholo}. We therefore distinguish
canopy cover, canopy height, and a canopy-volume proxy, while avoiding the claim
that any of them is measured biomass.

\subsection{Surface model, vegetation, and thermal mechanisms}

Surface-model information strengthens urban morphology by representing the combined
surface rather than footprints alone. Studies using LCZ, 2D/3D indicators,
SVF, and seasonal temperature consistently show that radiative openness and
vertical roughness contribute information beyond planimetric density
\citep{wang2018AssessinLocal,zhou2022UnderstaEffects,
ma2024XgboostBased,wang2025SeasonalEffects}. However, a surface model is not automatically
a building-height model. Terrain, vegetation, bridges, water artefacts, and
vertical datum uncertainty must be separated from built-form interpretation.
Our 5 m surface model is therefore used through surface-elevation range and standard
deviation together with independently derived floor-count height proxies.

Vegetation can affect temperature through shade, evapotranspiration, and
surface-energy partitioning, but effects depend on vegetation structure,
background climate, urban form, and season. Evidence spans phenology,
class-specific surface temperature, green-roof scenarios, and urban cooling
across nature-based-solution types \citep{kabano2021EvidenceUrban,
wang2023EffectivFactors,sanchez2025AssessinSpatial,wei2025UrbanCooling}.
Accordingly, vegetation enters both the classification ablation and the climate
response analysis, while causal cooling claims are excluded from cross-sectional
correlations.

\subsection{Deep learning and explainable urban-climate models}

Deep learning is useful when raster texture and figure-ground configuration
contain information that tabular summaries cannot preserve. LCZ research has
shown gains from residual CNNs, multilevel feature fusion, and combinations of
incomplete buildings with Sentinel imagery \citep{qiu2018FeatureImportan,
qiu2020MultilevFeature,yoo2020ImprovinLocal}. Multi-view urban-fabric models
extend this idea beyond a single overhead image and across spatial scales
\citep{fang2024UrbanclaDeep}. These advances motivate CLIMORFA's planned raster
and figure-ground branches.

Prediction alone is insufficient. XGBoost, CatBoost, SHAP, automated machine
learning, unsupervised clustering, and nonlinear models increasingly expose
thresholds, interactions, and geographically varying effects
\citep{ma2024XgboostBased,liu2025UnveilinDifferen,
xu2025CharacteThermal,wang2025RevealinImpact,
liu2026UnderstaNonlinea,wan2026AnalyzinImpact}. Interpretive frameworks also
emphasize mechanisms and seasonal differences rather than a ranking of feature
importance alone \citep{han2025PerspectUndersta,
wang2025SeasonalEffects}. We therefore report tree importance only as a first
explanation and require SHAP plus spatially explicit neural attribution before
any final deep-model claim. The general analogue is explainable urban-form
machine learning that connects model evidence to planning interpretation
\citep{karl2026RelationBetween}.

\subsection{Grid scale, network context, and spatial validation}

Scale is a substantive uncertainty. The support of SVF, urban form, green
pattern, and temperature changes the observed relationship
\citep{middel2018ViewFactor,wu2022ScaleEffect,hu2025ComparinGreen,
li2026InvestigEffect}. Urban mobility and heat exposure also operate along
networks rather than through Euclidean circles \citep{zhang2025MobilitiHeat}.
For green-space supply, a buffer counts nearby area but ignores network cost,
population competition, and shared destinations. We therefore use a network
2SFCA and retain 400/800/1,200 m thresholds as explicit sensitivity parameters.

Spatial validation is the companion requirement. Typology and climate studies
can exploit regional texture, city identity, or neighbourhood autocorrelation
unless folds are spatially separated. Comparing spatial blocks and held-out
districts provides a more demanding test of transfer. It also reveals when
apparently strong average performance hides local failure. Finally, climate
vulnerability cannot be inferred from morphology alone: heat risk, perceived
stress, and composite vulnerability combine exposure with social sensitivity
and adaptive capacity \citep{mashhoodi2024UrbanForm,iqbal2025DoesPerceive,
turner2025DevelopmValidati}. Our planning interpretation is therefore confined
to morphology and access signals and does not label vulnerable populations.


